\section{把代码审查设计成一个系统}

\subsection{审查质量首先取决于 PR 质量}

课程里一个非常实际的观点是：很多所谓``低质量审查''，根源并不在 reviewer，而在 pull request 本身过大、上下文不足、变更意图不清晰。要让审查有效，团队需要先把 PR 设计成可审查对象。

\begin{itemize}
    \item 控制单个 PR 的规模，让 reviewer 能在有限时间内建立完整心智模型
    \item 提供变更背景：问题是什么、为什么现在改、是否有替代方案
    \item 写清验证方法：跑了哪些测试、有哪些已知风险
    \item 对大改动做堆叠（stacked diffs）或按主题拆分，而不是一次性提交
\end{itemize}

\begin{importantbox}{审查效率不是 reviewer 一个人的 KPI}
如果作者提交的是一个 3000 行、跨十个模块、没有说明的 PR，那么再强的 reviewer 也只能给出表层反馈。高效审查是作者、工具和 reviewer 共同设计出来的流程能力。
\end{importantbox}

\subsection{审查意见的三层结构}

一条好的 review comment 通常可以拆成三层：

\begin{enumerate}
    \item \textbf{观察}：指出具体代码或行为
    \item \textbf{原因}：解释为什么这会造成风险、复杂度或歧义
    \item \textbf{建议}：给出替代方案、提问路径或后续改进方向
\end{enumerate}

\begin{knowledgebox}{从 ``挑错'' 到 ``共同设计''}
资深团队往往把 code review 看成设计讨论的最后一道关口，而不是上线前的抱怨区。评论最有价值的部分，通常不是指出 bug，而是帮助作者看到设计上的长期成本。
\end{knowledgebox}

\subsection{本章小结}

代码审查要想高质量，前提是 PR 本身可被审查。小而清晰的变更、明确的背景和验证信息，能把审查从机械 diff 浏览升级为真正的工程讨论。

